% Historical text from the preceding edition; retained without revision.
% This file is documentation, not a second manuscript. Citations use the paper bibliography.
\begin{abstract}
For a graph $H$, the three-color Ramsey number $R(H,H,H)$ is the least
$N$ such that every coloring of the edges of the complete graph $K_N$
with three colors contains a monochromatic copy of $H$. Dybizbański,
Dzido and Radziszowski stated the formula
$R(P_n,P_n,P_n)=2n-2+(n\bmod 2)$ as a conjecture in 2012, crediting
its suggestion to Faudree and Schelp (1975); $P_n$ is the path on $n$ vertices.
Gyárfás, Ruszinkó, Sárközy and Szemerédi proved this for all
sufficiently large $n$, and Dybizbański, Dzido and Radziszowski
determined $R(P_8,P_8,P_8)=14$ and $R(P_9,P_9,P_9)=17$, completing the
cases $n\le 9$ and leaving $n=10$ as the first open case. We prove
$R(P_{10},P_{10},P_{10})=18$: every $3$-coloring of the edges of $K_{18}$
contains a monochromatic path on $10$ vertices. The proof is a SAT
computation: symmetry breaking anchored on a monochromatic $P_9$, whose
existence is re-derived by the same method for the smaller cases, path
constraints generated during the search, and a cube-and-conquer split
into $15{,}956$ subproblems, each with an unsatisfiability certificate.
The argument is formalized in Lean~4, with certificates checked by
Lean's verified LRAT checker using trusted native evaluation. The default
\texttt{lake build} proves the value only with the certificate facts as
hypotheses. The unconditional theorem is built by the full check, which
regenerates every certificate and checks it in Lean ($249$ and $301$
CPU-hours in our two runs), or installs the certificates from a published
dataset, which skips the regeneration ($234$ and $284$ of those hours).
\end{abstract}

\paragraph{AI production.}
The proofs, the Lean formalization, the search code, the certificate checkers and the text were produced by AI models (from OpenAI and Anthropic), working under the author's direction.

\paragraph{Scope of verification.}
This is a computer-assisted proof: a Lean formalization with trusted
native evaluation. The statement, the reduction from colorings to
propositional formulas, the cover of the case split, and the lower bound
are proved in Lean (v4.32.2) from the definitions and checked by its
kernel. Two kinds of check are natively evaluated, through
\texttt{Lean.ofReduceBool}, which trusts Lean's compiler
(\texttt{Lean.trustCompiler}): the $15{,}964$ unsatisfiability certificate
checks (seven rungs, the cover and $15{,}956$ cubes), run by Lean's
verified LRAT checker, and the comparison identifying the parsed $K_{18}$
formula with the encoder's output (Section~\ref{sec:lean}). The
corresponding refutations, before their conversion to the compact form Lean
reads, were also checked independently of Lean by the verified checker
\texttt{cake\_lpr}, and a separate checker with its own code re-checked
every piece (Section~\ref{sec:separate}). No human expert has reviewed the
proof.

\section{Introduction}\label{sec:intro}

\subsection{The problem}
For graphs $G_1,G_2,G_3$, the Ramsey number $R(G_1,G_2,G_3)$ is the least
$N$ such that every coloring of the edges of $K_N$ with colors $1,2,3$
contains, for some $i$, a copy of $G_i$ all of whose edges have color $i$.
We write $\R(G)=R(G,G,G)$, and $P_n$ for the path on $n$ vertices (with
$n-1$ edges). Paths are not required to be induced.

Faudree and Schelp~\cite{FS75} studied $R(P_m,P_n,P_k)$ for paths of
different lengths and, ``from checking lower cases'', suggested a formula
for all of them~\cite[p.~159]{FS75}, whose diagonal case is the formula
below; Dybizbański, Dzido and Radziszowski~\cite{DDR16}
stated that case as a conjecture, and Radziszowski's dynamic
survey~\cite{Rad26} records it as the conjecture of Faudree and Schelp.

[The conjecture is retained verbatim in paper/paper.tex.]

The lower bound $\R(P_n)\ge 2n-2+(n\bmod2)$ holds for every $n$, by
Faudree and Schelp~\cite{FS75}: for odd $n$ by a blow-up of the
factorization of $K_4$ into three perfect matchings, and for even $n$ by
their general lower bound for paths of different lengths (proof of
Theorem~6). Gyárfás, Ruszinkó, Sárközy and Szemerédi~\cite{GRSS07} give
the blow-up for even $n$ as well (see also~\cite{DDR16} and
Section~\ref{sec:lower}). Figaj and
Łuczak~\cite{FL07} determined the three-color Ramsey numbers of long
even cycles asymptotically, which gives $R(P_{n_1},P_{n_2},P_{n_3})$
asymptotically as well~\cite[Section~6.4.1(a)]{Rad26}, and Gyárfás, Ruszinkó, Sárközy and
Szemerédi~\cite{GRSS07} proved Conjecture~\ref{conj:FS} for all
sufficiently large $n$, by the regularity method, with no explicit
threshold; Dybizbański, Dzido and Radziszowski report that their attempts
to extract a reasonable bound on the threshold from that proof
failed~\cite[Section~2]{DDR16}. Letzter~\cite{Let22} later gave a simpler
proof of the asymptotic bound $\R(P_n)\le(2+o(1))n$, with a sketch of a
corresponding stability result.

For small $n$ the conjecture is known for $n\le 9$. Table~\ref{tab:known}
lists the values with their sources as the survey~\cite[Section~6.4.1(c)]{Rad26}
attributes them, except for $n=5$ (see its caption). The values for $P_5$, $P_6$ and $P_7$ were obtained with
the help of computer algorithms~\cite[Section~2]{DDR16}. Dybizbański,
Dzido and Radziszowski~\cite{DDR16} (preprint 2012, published 2016) proved
$\R(P_8)=14$ and $\R(P_9)=17$ without computer help, which, in their words, ``leaves
$\R(P_{10})$ as the first open case''. They also attempted $P_{10}$ and
$P_{11}$ by the same method and did not complete these proofs; in their
words, ``between the first two open cases of Conjecture~2, namely the
questions whether it is true that $\R(P_{10})=18$ and $\R(P_{11})=21$, we
expect the former to be more difficult to prove'' (closing remarks
of~\cite{DDR16}). The latest revision of Radziszowski's survey
(revision \#18, April 24, 2026) still states that the conjecture ``is true
for $n\le 9$ \dots, and the first open case is that for
$P_{10}$''~\cite[Section~6.4]{Rad26}.

\begin{table}[ht]
\centering
\begin{tabular}{@{}lccccccc@{}}
\toprule
$n$ & 3 & 4 & 5 & 6 & 7 & 8 & 9\\
\midrule
$\R(P_n)$ & 5 & 6 & 9 & 10 & 13 & 14 & 17\\
source & \cite[Ea1]{Rad26} & \cite{Irv74} & \cite{YR94} & \cite{YR92} & \cite{Yan95} & \cite{DDR16} & \cite{DDR16}\\
\bottomrule
\end{tabular}
\medskip
\caption{The known values of $\R(P_n)$, with the sources the survey~\cite{Rad26}
gives, except for $n=5$: the survey cites~\cite{YR92}, which treats graphs
with five edges, while $P_5$, with four, is among the graphs
of~\cite{YR94}. For $n=3$ the survey notes that the value is easy to
obtain.}\label{tab:known}
\end{table}

For comparison, counting edges with the extremal theorem for paths gives
an upper bound. Erdős and Gallai~\cite{EG59} bounded the number of edges
of a graph with no long path; Faudree and Schelp~\cite{FS75}, and
independently Kopylov~\cite{Kop77}, determined it exactly (see
also~\cite[Theorem~4]{DDR16}). By their theorem, a $P_{10}$-free graph on $23$ vertices has at most
$2\binom92+\binom52=82$ edges, and $3\cdot 82<\binom{23}{2}$, so
$\R(P_{10})\le 23$.

\subsection{The result}

[The main theorem is retained verbatim in paper/paper.tex.]

The lower bound is the classical construction; the new part is the upper
bound $\R(P_{10})\le18$, which settles the case $n=10$ of
Conjecture~\ref{conj:FS}. The result is likely new: we found it neither
in the survey, which lists the case as the first open
one~\cite[Section~6.4]{Rad26}, nor in the literature we searched.

The prior exact results are those of Table~\ref{tab:known}: Irving's
$\R(P_4)=6$~\cite{Irv74}, Yang and Rowlinson's $\R(P_5)=9$~\cite{YR94} and
$\R(P_6)=10$~\cite{YR92} and Yang's $\R(P_7)=13$~\cite{Yan95}, the last
three obtained with computer help~\cite[Section~2]{DDR16}, and
Dybizbański, Dzido and Radziszowski's $\R(P_8)=14$ and
$\R(P_9)=17$~\cite{DDR16}, by hand. Each covers one value of $n$; none
bounds $\R(P_{10})$ from above by less than the counting bound $23$ below,
and the only other results bounding $\R(P_{10})$ that we found are
Galindez's computation $\R(C_{10})\le26$~\cite{Gal26} and Zhang's Lean
proof of a bound for all trees~\cite{ZhJ26}, each giving
$\R(P_{10})\le26$. On October
21, 2026 we rechecked the indexed forward citations of the canonical
papers: the one citation of~\cite{DDR16} that OpenAlex and Crossref
index, and the citations
of~\cite{GRSS07} since 2023 (among them~\cite{ABC+26}, on paths with few
color changes); none treats $\R(P_{10})$. The news is the value. The Lean
proof is how it is established and checked; we claim no novelty for the
method, whose ingredients have precedents~\cite{HKWB12,HS24,SNG+24,Sze26}.
The companion case $n=11$, $\R(P_{11})=21$, remains open; Dybizbański,
Dzido and Radziszowski showed that it would follow from the conjectured
value $\R(C_{10})=20$ for the $10$-cycle~\cite{DDR16}, which is itself the
first open case of the conjecture $\R(C_{2m})=4m$ for $m\ge3$, which
Dzido made in his 2005 thesis~\cite[Section~6.3.1(b)]{Rad26}; Benevides and
Skokan~\cite{BS09} proved it for large $m$.

The proof is a computation, built directly on the work of Dybizbański,
Dzido and Radziszowski. A monochromatic $P_9$ exists in every
$3$-coloring of $K_{18}$ because $\R(P_9)=17$~\cite{DDR16}, and fixing it
is the first and most effective reduction of the search. We do not take
that bound from the literature: a ladder of seven smaller computations,
$P_3$ on $K_5$ up to $P_9$ on $K_{17}$, re-derives the upper bounds
$\R(P_m)\le 5,6,9,10,13,14,17$ for $m=3,\dots,9$, each rung fixing the
monochromatic path that the rung below guarantees
(Section~\ref{sec:ladder}). The top rung, $P_{10}$ on $K_{18}$, is split
into $15{,}956$ subproblems (``cubes''); each is refuted by a SAT solver
with a proof certificate, and a separate certificate shows that the cubes
cover all cases (Section~\ref{sec:cnc}).

Lean~4~\cite{MU21} with Mathlib~\cite{Mathlib20} proves the reduction
from colorings to formulas: from the unsatisfiability of the rungs, the
cover formula and the cubes, it derives that every $3$-coloring of
$K_{18}$ has a monochromatic $P_{10}$. Every certificate (all $15{,}956$ cubes,
the cover and the seven rungs) is checked inside Lean by the verified LRAT
checker that ships with Lean (it underlies the \texttt{bv\_decide}
tactic~\cite{BBC+25}), run by native evaluation. The final theorem is the
Lean statement \texttt{ramsey3Path 10 = 18} (Section~\ref{sec:lean}).

\subsection{Related and concurrent work}\label{sec:related}
SAT solvers have settled several Ramsey-type questions, among them the
Boolean Pythagorean triples problem~\cite{HKM16}, by
cube-and-conquer~\cite{HKWB12} with a proof certificate, and the
three-color Ramsey number $R(4,3,3)=30$~\cite{CFIM16}, whose authors
expected certificates for it to be of unmanageable size; the empty hexagon
theorem~\cite{HS24} was later verified in Lean~\cite{SNG+24}. Szeider has
recently described a tool that imports LRAT certificates, including those
of cube-and-conquer computations, into Lean theorems through the same
verified checker of Lean's core library~\cite{Sze26}; our certificate
import (Section~\ref{sec:certs}) relies on the same checker, with its own
compact certificate encoding. On multicolor cycle
Ramsey numbers, Wesley~\cite{Wes25} used SAT solvers to prove
$R(C_3,C_6,C_6)=R(C_5,C_6,C_6)=15$ together with new bounds, and
Galindez~\cite{Gal26} has posted an exhaustive SAT computation, using SAT
Modulo Symmetries~\cite{KS24} and cube-and-conquer, giving
$\R(C_{10})\le 26$. That bound implies $\R(P_{10})\le 26$, since a
$10$-cycle contains a $P_{10}$. Shao, Jiang, Xu and Wang~\cite{SJXW26}
have posted a reproduction package for $R(C_4,C_6,C_8)=12$: after a hand
reduction in their paper, nine SAT instances, each refuted with a RUP
certificate. A public SAT project on three-color cycle
Ramsey numbers~\cite{Zha26} is also active. For trees, Zhang~\cite{ZhJ26}
has posted a Lean proof that $R_k(T)\le k(n-2)+2$ for every tree $T$ on
$n\ge2$ vertices, as a corollary of a machine-checked proof of the
Erdős--Sós conjecture found by GPT-6 Astra and packaged by Adamczewski; for three colors and $P_{10}$
it gives $26$. Three-color Ramsey numbers of
paths of different lengths, and of paths against cycles, have been
determined in many further cases~\cite{Dzi05,DF09,OR11,MORS15}, and
Radziszowski's survey~\cite{Rad26} records many further computer-assisted
results on small Ramsey numbers.

\subsection{Organization}
Section~\ref{sec:lower} gives the lower bound. Section~\ref{sec:encoding}
describes the propositional encoding and the two symmetry reductions,
Section~\ref{sec:ladder} the ladder, and Section~\ref{sec:cnc} the case
split on $K_{18}$ and why it is complete. Section~\ref{sec:lean} describes
the Lean formalization and exactly what it trusts;
Section~\ref{sec:separate} the checks made independently of Lean, and
which steps of that route are machine-checked and which are argued.
Section~\ref{sec:replay} gives the computing cost and how to replay the
proof.

